\documentclass[11pt,a4paper]{article}
\usepackage[T1]{fontenc}
\usepackage[utf8]{inputenc}
\usepackage{amsmath,amssymb,amsthm}
\usepackage[margin=1in]{geometry}
\usepackage[colorlinks=true,linkcolor=blue,urlcolor=blue]{hyperref}
\theoremstyle{plain}
\newtheorem{theorem}{Theorem}
\newtheorem{lemma}[theorem]{Lemma}
\newtheorem{proposition}[theorem]{Proposition}
\newtheorem{corollary}[theorem]{Corollary}
\theoremstyle{definition}
\newtheorem{remark}[theorem]{Remark}

\title{The white squares of a board: OEIS A230704}
\author{Adrian Perez Fontelles\\ \small Independent researcher}
\date{7 September 2026}
\begin{document}
\maketitle

\begin{abstract}
OEIS A230704 carries an empirical recurrence of order $50$ contributed by R.~H.~Hardin. It is
true, and it is decidable rather than empirical. The entry's condition is decided by a bounded amount of state carried down the array; the
admissible windows are the vertices of a finite digraph and the arrays counted are exactly the
walks in it. Hence $a(n)$ is a walk count on $S=166$ vertices and is $C$-finite, and whether the
conjectured recurrence annihilates it is settled by a finite exact computation, with the
Cayley--Hamilton theorem bounding the work.
\end{abstract}

\noindent\small 2020 Mathematics Subject Classification. 05A15, 05B45, 15B36.\normalsize

\section{The conjecture}

OEIS A230704 is ``Number of (n+3)X(4+3) 0..2 white square subarrays x(i,j) with each element diagonally or antidiagonally next to at least one element with value (x(i,j)+1) mod 3, and upper left element zero.''. It has offset $1$ and begins
\[
102, 800, 16536, 121382, 2382398, 17497342, 342031378, 2508344646,\ \dots
\]
The entry states, as a conjecture contributed by its author and never marked settled:
\begin{quote}
Empirical: a(n) = 185*a(n-2) -6851*a(n-4) +145707*a(n-6) -2849231*a(n-8) +30720183*a(n-10) -36007996*a(n-12) -743385180*a(n-14) +287330815*a(n-16) +8627923017*a(n-18) -12118273932*a(n-20) -29408445279*a(n-22) +45675171837*a(n-24) -13807146112*a(n-26) -2411662568755*a(n-28) +3802259458200*a(n-30) +3456476952882*a(n-32) +4249525502473*a(n-34) -5769431443156*a(n-36) -7037014662528*a(n-38) +2632063276930*a(n-40) +18416679443602*a(n-42) +9430615399362*a(n-44) +259673859628*a(n-46) +64576080600*a(n-48) +325261872*a(n-50)
\end{quote}
As of the ``Last modified'' line on the live entry (Jul 23 2025 05:57:44, revision 6) this is still
recorded as empirical, and nothing on the entry records it as proved.

\section{The white squares are a board of their own}

The board has $7$ columns and grows one row at a time, but only the cells with $i+j$ even carry a value, from the alphabet $\{0,1,2\}$. The diagonal and antidiagonal neighbours of such a cell again have $i+j$ even, so the white squares form a board of their own and the entry's condition never refers to anything else: each white cell must have at least one neighbour, diagonally or antidiagonally, holding the value $x+1$ modulo $3$ where $x$ is its own.

That looks one row up and one row down, so three consecutive rows settle the condition on the middle one. Carry two rows as the state, judge a row when the row after it arrives, and judge the row with nothing after it by the end vector. A cell at the edge of the board simply has fewer neighbours, and having none of the right value is a failure rather than a special case.

The white cells of a row occupy alternate columns, and which alternate ones depends on the parity of that row, so the two shapes alternate down the board; the state carries that parity in the shape of the row it holds.

The entry also fixes the top left cell at $0$, which restricts the states an array may begin in and nothing else.

Merging vertices with identical futures leaves $S=166$. An admissible board is exactly a walk, so with $M$ the adjacency matrix and $\iota,\tau$ the vectors recording which states may begin and end a board,
\[
a(n)\;=\;\iota^{\!\top}M^{\,n+2}\tau ,
\]
the exponent fixed by the entry's own shape and checked against its published terms. In
particular $a$ is $C$-finite.

\section{The proof}

\begin{lemma}\label{lem:crit}
Let $q(t)=t^{50}-\sum_i c_it^{50-i}$ be the characteristic polynomial of the
conjectured recurrence. The residual $a(n)-\sum_ic_ia(n-i)$ equals
$\iota^{\!\top}M^{\,j}q(M)\tau$ for the corresponding walk index $j$, so the recurrence
holds from some point on if and only if $\iota^{\!\top}M^{j}q(M)\tau=0$ for all large $j$,
and it is enough to examine $j<S$.
\end{lemma}

\begin{proof}
Factor $M^{j}$ out of $M^{j+50}-\sum_ic_iM^{j+50-i}$. For the bound, $M^{S}$
is an integer combination of $I,M,\dots,M^{S-1}$ by Cayley--Hamilton, so the numbers
$u_j=\iota^{\!\top}M^{j}q(M)\tau$ obey the monic recurrence given by the characteristic
polynomial of $M$; $S$ consecutive zeros force every later one, and the last nonzero $u_j$
pins the threshold exactly. A constant factor in front of $a$ divides out of the residual and
changes nothing here.
\end{proof}

\begin{theorem}
\[
a(n)\;=\;185\,a(n-2) - 6851\,a(n-4) + 145707\,a(n-6) - 2849231\,a(n-8) + 30720183\,a(n-10) - 36007996\,a(n-12) - 743385180\,a(n-14) + 287330815\,a(n-16) + 8627923017\,a(n-18) - 12118273932\,a(n-20) - 29408445279\,a(n-22) + 45675171837\,a(n-24) - 13807146112\,a(n-26) - 2411662568755\,a(n-28) + 3802259458200\,a(n-30) + 3456476952882\,a(n-32) + 4249525502473\,a(n-34) - 5769431443156\,a(n-36) - 7037014662528\,a(n-38) + 2632063276930\,a(n-40) + 18416679443602\,a(n-42) + 9430615399362\,a(n-44) + 259673859628\,a(n-46) + 64576080600\,a(n-48) + 325261872\,a(n-50)
\]
for every $n>48$.
\end{theorem}

\begin{proof}
The vector $w=q(M)\tau$ was formed by $50$ matrix--vector products in exact integer
arithmetic and $\iota^{\!\top}M^{j}w$ evaluated for $j=0,1,\dots$, again exactly; those
integers vanish from the index corresponding to $n=48+1$ onwards and the run continued
until the working vector was identically zero, which settles every larger $j$ at once.
\end{proof}

\section{Verification}

Four checks, each independent of the algebra above.

First, the model reproduces all $16$ terms the entry publishes, exactly, in integer
arithmetic. This is what ties the model to the entry: the digraph is built by reading the
entry's English, and a misreading gives different counts.

Second, the reading was pinned before the digraph was written, by a program that enumerates
every filling of the white cells one by one and tests, for each of them, whether some
diagonal or antidiagonal neighbour holds the required value. That program shares no code with the transfer matrix, so an error in one does not
hide an error in the other.

Third, the conjectured recurrence was evaluated directly on the published terms, with no
matrices involved, and holds wherever the proved range and the data overlap.

Fourth, the annihilation test of Lemma~\ref{lem:crit} was carried out in exact integer
arithmetic on the whole vector rather than on a sampled prefix.

\begin{thebibliography}{9}
\bibitem{oeis} The OEIS Foundation, \emph{The On-Line Encyclopedia of Integer
Sequences}, \url{https://oeis.org}, 2026. Sequence A230704.
\bibitem{stanley} R.~P.~Stanley, \emph{Enumerative Combinatorics}, Volume 1, 2nd ed.,
Cambridge University Press, 2011. (Transfer-matrix method, Section 4.7.)
\bibitem{fs} P.~Flajolet and R.~Sedgewick, \emph{Analytic Combinatorics}, Cambridge
University Press, 2009.
\bibitem{horn} R.~A.~Horn and C.~R.~Johnson, \emph{Matrix Analysis}, 2nd ed., Cambridge
University Press, 2013.
\end{thebibliography}

\end{document}
